\chapter{Análisis del problema, requisitos, planificación y presupuesto}
\label{chap:metodologia-plan}

\lettrine{L}{a} ejecución del TFG se plantea como un desarrollo incremental de arquitectura, orientado a obtener un sistema funcional por fases y a reducir riesgo técnico en integración.

\section{Descripción general del sistema a desarrollar}

El sistema se define como una arquitectura ciberfísica para captura y gestión de datos inerciales.
La propuesta integra adquisición, comunicación, persistencia y visualización/control bajo una estructura modular, de forma que Android y el nodo embebido puedan interoperar sin depender de implementaciones rígidas.

\section{Requisitos funcionales}

Los requisitos funcionales principales del sistema son:

\begin{itemize}
  \item capturar datos inerciales desde nodo móvil Android;
  \item capturar o descargar sesiones desde nodo embebido ESP32+IMU;
  \item almacenar sesiones con metadatos trazables;
  \item visualizar estado de nodo, sesiones y resultados de forma operativa;
  \item exportar sesiones para análisis externo;
  \item soportar comandos básicos de control remoto (inicio, parada, estado y configuración).
\end{itemize}

\section{Requisitos no funcionales}

Además de la funcionalidad visible, el diseño debe cumplir:

\begin{itemize}
  \item modularidad entre adquisición, comunicación, backend y visualización;
  \item mantenibilidad para evolucionar backend/transporte sin rehacer el sistema;
  \item bajo coste de despliegue en entorno de TFG;
  \item trazabilidad de sesiones y de parámetros de ejecución;
  \item facilidad de operación en laboratorio y escenarios remotos.
\end{itemize}

\section{Restricciones técnicas}

Las restricciones más relevantes del proyecto son:

\begin{itemize}
  \item heterogeneidad de sensores y temporización en Android;
  \item limitaciones de recursos y estabilidad de integración en ESP32;
  \item necesidad de definir contratos de datos comunes entre nodos;
  \item tiempo de desarrollo limitado para cerrar integración extremo a extremo;
  \item incertidumbre sobre qué transporte remoto aporta mejor equilibrio en este alcance.
\end{itemize}

\section{Casos de uso principales}

Los casos de uso que estructuran el sistema son:

\begin{itemize}
  \item captura y registro local desde Android;
  \item operación de captura desde ESP32 con supervisión desde la app;
  \item consulta histórica y reprocesado de sesiones;
  \item comparación de sesiones entre nodos/dispositivos;
  \item configuración remota y consulta de estado.
\end{itemize}

\figplaceholder{%
Diagrama de casos de uso recomendado:\\
usuario, nodo móvil, nodo embebido, backend, panel de consulta
}{Casos de uso principales del sistema.}
\label{fig:placeholder-casos-uso}

\section{Enfoque de trabajo}

Se adopta una \textbf{metodología incremental e iterativa}.
La razón principal es separar riesgos por capas: primero análisis del problema y diseño arquitectónico, después integración entre nodos y finalmente validación operativa.

Este enfoque permite validar decisiones con entregables funcionales y evitar bloqueos por dependencias tempranas de infraestructura. El trabajo parte de un análisis inicial del problema, del estado del arte y de las tecnologías disponibles para definir la arquitectura final del sistema.

De forma resumida, la estrategia por bloques funcionales es:

\begin{enumerate}
  \item captura y persistencia de sesión en local;
  \item interfaz operativa Android y gestión de sesiones;
  \item integración local con ESP32 vía interfaz de red;
  \item incorporación progresiva de supervisión y acceso externo a datos;
  \item evaluación final del sistema y propuestas de mejora.
\end{enumerate}

\section{Fases del proyecto}

\begin{itemize}
  \item \textbf{Fase 1.} Análisis de trabajos relacionados en sensores inerciales, análisis de movimiento y arquitecturas de captura de datos.
  \item \textbf{Fase 2.} Estudio de limitaciones y características de sensores IMU y técnicas de análisis aplicables.
  \item \textbf{Fase 3.} Definición de la arquitectura del sistema y de la interacción entre módulos.
  \item \textbf{Fase 4.} Implementación de captura y gestión de datos en diferentes plataformas.
  \item \textbf{Fase 5.} Evaluación del sistema desarrollado y análisis de comportamiento, limitaciones y mejoras.
  \item \textbf{Fase 6.} Aplicación del sistema al caso práctico propuesto.
  \item \textbf{Fase 7.} Documentación final y presentación.
\end{itemize}

\section{Planificación de operación remota (borrador)}
\label{sec:plan-operacion-remota}

La operación remota se plantea como evolución del sistema y no como requisito cerrado desde el inicio.
En el estado actual, el objetivo es dejar una base operativa con sincronización de sesiones y estado de nodos.

Para ello se propone:

\begin{itemize}
  \item \textbf{Implementación base:} integración con mecanismo de almacenamiento/sincronización externo según viabilidad del entorno de pruebas.
  \item \textbf{Sincronización mínima:} metadatos de sesión, estado de nodo y referencias a ficheros.
  \item \textbf{Control básico:} comandos start/stop/set\_config/status bajo contrato común.
  \item \textbf{Evolución opcional:} evaluar alternativas de comunicación si mejoran latencia, robustez o escalabilidad.
\end{itemize}

Este planteamiento evita comprometer decisiones de infraestructura que no estén justificadas por evidencia experimental.

\section{Plan de pruebas y control de riesgos}

Los riesgos principales y su mitigación son:

\begin{itemize}
  \item \textbf{Desalineación entre nodos:} fijar contrato de sesión y validarlo en ambos modos.
  \item \textbf{Inestabilidad de integración embebida:} mantener diagnósticos de conexión, bus y captura.
  \item \textbf{Acoplamiento excesivo al backend:} desacoplar capa de comunicación y persistencia.
  \item \textbf{Sobrecarga de alcance:} priorizar MVP funcional antes de añadir extensiones opcionales.
\end{itemize}

\section{Material y medios necesarios}

Los medios necesarios para ejecutar el proyecto son:

\begin{itemize}
  \item PC para desarrollo, pruebas y análisis de resultados.
  \item Dispositivo móvil Android con sensores inerciales.
  \item Plataforma embebida para captura de datos con conectividad inalámbrica.
  \item Sensor IMU externo (acelerómetro y/o giróscopo).
  \item Herramientas de desarrollo software y entornos de programación.
  \item Herramientas para almacenamiento, visualización y análisis de datos.
\end{itemize}

\section{Presupuesto orientativo}

El trabajo prioriza bajo coste de hardware:

\begin{itemize}
  \item teléfono Android como nodo principal;
  \item placa ESP32 disponible para integración embebida;
  \item IMU MPU6050 para pruebas de adquisición dedicada.
\end{itemize}

El coste principal se concentra en tiempo de desarrollo, integración y validación del sistema extremo a extremo.
